\ifdefined\gkver\else\def\gkver{student}\fi
\documentclass[\gkver]{gaokaozhenti}
\begin{document}

\chapter{2017年天津理综}
%% paperType: 理综物理部分

\begin{enumerate}
\item
%% number: 1
%% paperName: 2017年天津理综第 1 题 6 分
%% typeId: 1
%% type: 单选题
%% chapter: 原子和原子核
%% point: 核聚变
%% method: 无
%% score: 6
%% degree: 950
%% duplicateId: 0
%% body:
我国自主研发制造的国际热核聚变核心部件在国际上率先通过权威机构认证，这是我国对国际热核聚变项目的重大贡献。下列核反应方程中属于聚变反应的是\xzanswer{A}

\onepicture[h=3.0cm]{figs/01.jpg}

\fourchoices[answer=A]
{$\ce{^{2}_{1}H + ^{3}_{1}H -> ^{4}_{2}He + ^{1}_{0}n}$}
{$\ce{^{14}_{7}N + ^{4}_{2}He -> ^{17}_{8}O + ^{1}_{1}H}$}
{$\ce{^{4}_{2}He + ^{27}_{13}Al -> ^{30}_{15}P + ^{1}_{0}n}$}
{$\ce{^{235}_{92}U + ^{1}_{0}n -> ^{144}_{56}Ba + ^{89}_{36}Kr + 3 ^{1}_{0}n}$}

%% answer: A
%% memo:
\memoanswer{
两个轻核结合成质量较大的核的反应是核聚变反应，故 A 正确，BCD 错误。
}

\item
%% number: 2
%% paperName: 2017年天津理综第 2 题 6 分
%% typeId: 1
%% type: 单选题
%% chapter: 光学
%% point: 光的折射
%% method: 无
%% score: 6
%% degree: 550
%% duplicateId: 0
%% body:
明代学者方以智在《阳燧倒影》中记载：“凡宝石面凸，则光成一条，有数棱则必有一面五色”，表明白光通过多棱晶体折射会发生色散现象。如图所示，一束复色光通过三棱镜后分解成两束单色光 a、b，下列说法正确的是\xzanswer{D}

\onepicture[h=3.0cm]{figs/02.png}

\fourchoices[answer=D]
{若增大入射角 $i$，则 b 光先消失}
{在该三棱镜中 a 光波长小于 b 光}
{a 光能发生偏振现象，b 光不能发生}
{若 a、b 光分别照射同一光电管都能发生光电效应，则 a 光的遏止电压低}

%% answer: D
%% memo:
\memoanswer{
由折射定律及几何关系可知，若增大入射角 $i$，在另一界面，a、b 光的入射角都将减小，都不会发生全反射，两光都不会消失，故 A 错误；由题图及光的折射定律可知棱镜对 b 光的折射率较大，由折射率与频率的关系可知 b 光频率较大，由光速与折射率关系可知在棱镜中 b 光速度较小，由 $v=\lambda f$ 可知在棱镜中 b 光波长较小，故 B 错误；光是横波，都能发生偏振现象，故 C 错误；由光电效应方程可知，a、b 两光分别照射同一光电管发生光电效应时，被 a 光光子激发出的光电子的最大初动能较小，由 $E_{km}=eU_{c}$ 可知 a 光的遏止电压较低，故 D 正确。
}

\item
%% number: 3
%% paperName: 2017年天津理综第 3 题 6 分
%% typeId: 1
%% type: 单选题
%% chapter: 电磁感应
%% point: 电磁感应受力分析
%% method: 无
%% score: 6
%% degree: 650
%% duplicateId: 0
%% body:
如图所示，两根平行金属导轨置于水平面内，导轨之间接有电阻 $R$。金属棒 ab 与两导轨垂直并保持良好接触，整个装置放在匀强磁场中，磁场方向垂直于导轨平面向下。现使磁感应强度随时间均匀减小，ab 始终保持静止，下列说法正确的是\xzanswer{D}

\onepicture[h=3.4cm]{figs/03.png}

\fourchoices[answer=D]
{ab 中的感应电流方向由 b 到 a}
{ab 中的感应电流逐渐减小}
{ab 所受的安培力保持不变}
{ab 所受的静摩擦力逐渐减小}

%% answer: D
%% memo:
\memoanswer{
由楞次定律可知，ab 中的感应电流方向由 a 到 b，故 A 错误；由于闭合电路面积不变，磁感应强度均匀减小，穿过闭合电路的磁通量均匀减小，由法拉第电磁感应定律 $E=\frac{\Delta \phi}{\Delta t}$ 可知，感应电动势恒定不变，由欧姆定律可知 ab 中的感应电流恒定不变，故 B 错误；由于磁感应强度减小，由安培力公式 $F=BIL$ 可知，ab 所受安培力逐渐减小，故 C 错误；由共点力平衡条件可知，ab 所受静摩擦力逐渐减小，故 D 正确。
}

\item
%% number: 4
%% paperName: 2017年天津理综第 4 题 6 分
%% typeId: 1
%% type: 单选题
%% chapter: 曲线运动
%% point: 圆周运动的应用
%% method: 无
%% score: 6
%% degree: 650
%% duplicateId: 0
%% body:
“天津之眼”是一座跨河建设、桥轮合一的摩天轮，是天津市的地标之一。摩天轮悬挂透明座舱，乘客随座舱在竖直面内做匀速圆周运动。下列叙述正确的是\xzanswer{D}

\onepicture[h=4.0cm]{figs/04.png}

\fourchoices[answer=D]
{摩天轮转动过程中，乘客的机械能保持不变}
{在最高点，乘客重力大于座椅对他的支持力}
{摩天轮转动一周的过程中，乘客重力的冲量为零}
{摩天轮转动过程中，乘客重力的瞬时功率保持不变}

%% answer: D
%% memo:
\memoanswer{
摩天轮转动过程中，乘客动能不变，重力势能变化，则机械能变化，故 A 错误；在最高点，乘客有竖直向下的向心加速度，由牛顿第二定律可知，乘客的重力大于座椅对他的支持力，故 B 正确；摩天轮转动一周的过程中，所受重力恒定不变，重力的冲量不等于零，故 C 错误；转动过程中，乘客速度方向与重力方向的夹角变化，重力的瞬时功率发生变化，故 D 错误。
}

\item
%% number: 5
%% paperName: 2017年天津理综第 5 题 6 分
%% typeId: 1
%% type: 单选题
%% chapter: 机械波
%% point: 波的图象
%% method: 无
%% score: 6
%% degree: 650
%% duplicateId: 0
%% body:
手持较长软绳端点 O 以周期 $T$ 在竖直方向上做简谐运动，带动绳上的其他质点振动形成简谐波沿绳水平传播，示意如图。绳上有另一质点 P，且 O、P 的平衡位置间距为 $L$。$t=0$ 时，O 位于最高点，P 的位移恰好为零，速度方向竖直向上，下列判断正确的是\xzanswer{C}

\onepicture[h=2.6cm]{figs/05.png}

\fourchoices[answer=C]
{该简谐波是纵波}
{该简谐波的最大波长为 $2L$}
{$t=\frac{T}{8}$ 时，P 在平衡位置上方}
{$t=\frac{3T}{8}$ 时，P 的速度方向竖直向上}

%% answer: C
%% memo:
\memoanswer{
波的传播方向与绳上各质点的振动方向垂直，是横波，故 A 错误；由题意可得 $L=\left(n+\frac{1}{4}\right)\lambda$，解得最大波长为 $4L$，故 B 错误；由于 $t=0$ 时刻质点 P 过平衡位置向上运动，则 $t=\frac{T}{8}$ 时刻 P 正向上运动且未到波峰，在平衡位置上方，故 C 正确；由于 $\frac{1}{2}T>\frac{3T}{8}>\frac{1}{4}T$，$t=\frac{3T}{8}$ 时 P 的速度方向竖直向下，故 D 错误。
}

\item
%% number: 6
%% paperName: 2017年天津理综第 6 题 6 分
%% typeId: 2
%% type: 多选题
%% chapter: 交流电
%% point: 交流电
%% method: 无
%% score: 6
%% degree: 650
%% duplicateId: 0
%% body:
在匀强磁场中，一个 100 匝的闭合矩形金属线圈，绕与磁感线垂直的固定轴匀速转动，穿过该线圈的磁通量随时间按图示正弦规律变化。设线圈总电阻为 $2 \UO$，则\xzanswer{AD}

\onepicture[h=3.6cm]{figs/06.png}

\fourchoices[answer=AD]
{$t=0$ 时，线圈平面平行于磁感线}
{$t=1 \Us$ 时，线圈中的电流改变方向}
{$t=1.5 \Us$ 时，线圈中的感应电动势最大}
{一个周期内，线圈产生的热量为 $8\pi^{2} \UJ$}

%% answer: AD
%% memo:
\memoanswer{
$\Phi-t$ 图象的斜率表示磁通量的变化率，由图象可知，$t=0$ 时刻穿过线圈的磁通量为零，线圈平面平行于磁感线，故 A 正确；由图象斜率意义及法拉第电磁感应定律可知，$t=0.5 \Us$ 或 $t=1.5 \Us$ 时线圈中的感应电动势为零，线圈中的电流改变方向，故 BC 错误；由法拉第电磁感应定律 $e=n\frac{\Delta \phi}{\Delta t}$ 及正弦函数的导数可知，感应电动势的最大值 $E_{\mathrm{m}}=n\times \Phi_{\mathrm{m}}\times \frac{2\pi}{T}=4\pi \UV$，其有效值 $E=\frac{E_{\mathrm{m}}}{\sqrt{2}}$，则一个周期内线圈产生的热量 $Q=\frac{E^{2}}{R}T=8\pi^{2} \UJ$，故 D 正确。
}

\item
%% number: 7
%% paperName: 2017年天津理综第 7 题 6 分
%% typeId: 2
%% type: 多选题
%% chapter: 电场
%% point: 带电粒子的轨迹类问题
%% method: 无
%% score: 6
%% degree: 550
%% duplicateId: 0
%% body:
如图所示，在点电荷 $Q$ 产生的电场中，实线 MN 是一条方向未标出的电场线，虚线 AB 是一个电子只在静电力作用下的运动轨迹。设电子在 A、B 两点的加速度大小分别为 $a_{A}$、$a_{B}$，电势能分别为 $E_{pA}$、$E_{pB}$。下列说法正确的是\xzanswer{BC}

\onepicture[h=3.4cm]{\includesvg{figs/07}}

\fourchoices[answer=BC]
{电子一定从 A 向 B 运动}
{若 $a_{A}>a_{B}$，则 $Q$ 靠近 M 端且为正电荷}
{无论 $Q$ 为正电荷还是负电荷，一定有 $E_{pA}<E_{pB}$}
{B 点电势可能高于 A 点电势}

%% answer: BC
%% memo:
\memoanswer{
由于电子运动情况未知，无法判断电子的运动方向，故 A 错误；由电子运动轨迹及曲线运动的特点可知，电子在电场中所受电场力方向向左，则电场线 MN 上的电场方向向右，B 点电势低于 A 点电势，故 D 错误；由于 B 点电势低于 A 点电势，因此，电子在 A 点的电势能小于在 B 点的电势能，故 C 正确；若 $a_{A}>a_{B}$，则电子在 A 端时所受电场力较大，离 Q 较近，则 Q 靠近 M 端，由 B 点电势低于 A 点电势可得 Q 带正电，故 B 正确。
}

\item
%% number: 8
%% paperName: 2017年天津理综第 8 题 6 分
%% typeId: 2
%% type: 多选题
%% chapter: 力物体的平衡
%% point: 动态平衡
%% method: 无
%% score: 6
%% degree: 450
%% duplicateId: 0
%% body:
如图所示，轻质不可伸长的晾衣绳两端分别固定在竖直杆 M、N 上的 a、b 两点，悬挂衣服的衣架钩是光滑的，挂于绳上处于静止状态。如果只人为改变一个条件，当衣架静止时，下列说法正确的是\xzanswer{AB}

\onepicture[h=4.2cm]{figs/08.png}

\fourchoices[answer=AB]
{绳的右端上移到 $b'$，绳子拉力不变}
{将杆 N 向右移一些，绳子拉力变大}
{绳的两端高度差越小，绳子拉力越小}
{若换挂质量更大的衣服，则衣服架悬挂点右移}

%% answer: AB
%% memo:
\memoanswer{
由于挂钩光滑，对挂钩由共点力平衡条件可知，平衡时两侧绳与竖直方向的夹角相等，两侧绳的拉力相等，有 $G=2F\cos\theta$。由几何关系可知，绳的右端上移到 $b'$ 重新平衡后，挂钩两侧绳与竖直方向的夹角未变，由共点力平衡条件可知，绳子的拉力不变，故 A 正确；将杆 N 向右移一些重新平衡后，挂钩两侧绳与竖直方向的夹角变大，由 $G=2F\cos\theta$ 可知绳的拉力变大，故 B 正确；由选项 A 分析可得，绳的两端高度差改变，绳子拉力可能不变，故 C 错误；同理可知，若换挂质量更大的衣服，绳的拉力变大，但挂钩位置不动，故 D 错误。
}

\item
%% number: 9
%% paperName: 2017年天津理综第 9 题 3 分
%% typeId: 4
%% type: 实验题
%% chapter: 电路
%% point: 测电动势与内阻
%% method: 无
%% score: 3
%% degree: 350
%% duplicateId: 0
%% body:
（1）我国自主研制的首艘货运飞船“天舟一号”发射升空后，与已经在轨运行的“天宫二号”成功对接形成组合体。假设组合体在距地面高度为 $h$ 的圆形轨道上绕地球做匀速圆周运动，已知地球半径为 $R$，地球表面重力加速度为 $g$，且不考虑地球自转的影响。则组合体运动的线速度大小为\tkanswer[2.4cm]{$R\sqrt{\frac{g}{R+h}}$}，向心加速度大小为\tkanswer[2.6cm]{$\frac{R^{2}}{(R+h)^{2}}g$}。
\onepicture[h=3.0cm, align=right]{figs/09a.png}

（2）如图所示，打点计时器固定在铁架台上，使重物带动纸带从静止开始自由下落，利用此装置验证机械能守恒定律。
\onepicture[h=4.2cm, align=right]{figs/09b.png}
\begin{steps}
	\item
	对于该实验，下列操作中对减小实验误差有利的是\tkanswer{AB}。

	（A）重物选用质量和密度较大的金属锤

	（B）两限位孔在同一竖直面内上下对正

	（C）精确测量出重物的质量

	（D）用手托稳重物，接通电源后，撒手释放重物
	\item
	某实验小组利用上述装置将打点计时器接到 $50 \UHz$ 的交流电源上，按正确操作得到了一条完整的纸带，由于纸带较长，图中有部分未画出，如图所示。纸带上各点是打点计时器打出的计时点，其中 O 点为纸带上打出的第一个点。重物下落高度应从纸带上计时点间的距离直接测出，利用下列测量值能完成验证机械能守恒定律的选项有\tkanswer{BC}。
	\onepicture[w=12cm, align=right]{figs/09c.png}

	（A）OA、AD 和 EG 的长度

	（B）OC、BC 和 CD 的长度

	（C）BD、CF 和 EG 的长度

	（D）AC、BD 和 EG 的长度
\end{steps}

（3）某探究性学习小组利用如图 \ref{2017天津理综09d} 所示的电路测量电池的电动势和内阻。其中电流表 A$_{1}$ 的内阻 $r_{1}=1.0 \UkO$，电阻 $R_{1}=9.0 \UkO$，为了方便读数和作图，给电池串联一个 $R_{0}=3.0 \UO$ 的电阻。
\twopicture[labelA=2017天津理综09d, labelB=2017天津理综09e, align=right, hA=3.2cm, hB=3.0cm]
{figs/09d.jpg}
{figs/09e.jpg}
\begin{steps}
	\item
	按图示电路进行连接后，发现 aa$'$、bb$'$ 和 cc$'$ 三条导线中，混进了一条内部断开的导线。为了确定哪一条导线内部是断开的，将电键 S 闭合，用多用电表的电压挡先测量 a、b$'$ 间电压，读数不为零，再测量 a、a$'$ 间电压，若读数不为零，则一定是\tkanswer{aa$'$}导线断开；若读数为零，则一定是\tkanswer{bb$'$}导线断开。
	\item
	排除故障后，该小组顺利完成实验。通过多次改变滑动变阻器触头位置，得到电流表 A$_{1}$ 和 A$_{2}$ 的多组 $I_{1}$、$I_{2}$ 数据，作出图象如图 \ref{2017天津理综09e} 所示。由 $I_{1}-I_{2}$ 图象得到的电池的电动势 $E=$ \tkanswer{1.4} $\UV$，内阻 $r=$ \tkanswer{0.5} $\UO$。
\end{steps}

%% answer: 
\jdanswer{
\begin{enumerate}
	\item
	$R\sqrt{\frac{g}{R+h}}$，$\frac{R^{2}}{(R+h)^{2}}g$

	\item
	AB，BC

	\item
	aa$'$，bb$'$；1.4（1.36~1.44 均可），0.5（0.4~0.6 均可）

\end{enumerate}
}
%% memo:
\memoanswer{
（1）对组合体环绕地球的运动，由牛顿第二定律及万有引力定律得 $G\frac{Mm}{(R+h)^{2}}=\frac{mv^{2}}{R+h}=ma$，在地面有 $g=G\frac{M}{R^{2}}$。解得组合体运动的向心加速度大小 $a=\frac{R^{2}}{(R+h)^{2}}g$，线速度大小 $v=R\sqrt{\frac{g}{R+h}}$。

（2）①选用质量和密度较大的金属锤，空气的阻力可忽略不计，有利于减小实验误差，故 A 正确；两限位孔在同一竖直平面内且上下对正，可减小纸带所受阻力，选项 B 正确；实验中应先接通电源再释放重物，是正确的实验步骤，与误差无关，本实验要验证的是 $gh=\frac{1}{2}v^{2}$，实验中不必测量出重物的质量，故 CD 错误。

②由 OC 的长度可计算出从打 O 点到打 C 点重物减小的重力势能，由 BC、CD 的长度及交流电的频率可计算出打 C 点时重物的速度，从而计算出从打 O 点到打 C 点重物增加的动能，故 B 正确；由 CF 的长度可计算出从打 C 点到打 F 点重物减小的重力势能，由 BD、EG 的长度及交流电的频率可计算出打 C、F 点时重物的速度，从而计算出从打 C 点到打 F 点重物增加的动能，故 C 正确；同理可得 AD 不能得出结论，故 AD 错误。

（3）①测量 a、b$'$ 间的电压时读数不为零，说明从 b$'$ 到电源正极是连通的，从 a 到电源负极是连通的，若测 a、a$'$ 间的电压时读数不为零，说明从 a$'$ 到电源正极是连通的，从 a 到电源负极是连通的，则断开的一定是导线 aa$'$；若测 a、a$'$ 间的电压时读数为零，说明 a$'$ 与电源正极未连通，一定是导线 bb$'$ 断开。

②对实验电路，由电路欧姆定律得

$E=I_{1}(r_{1}+R_{1})+(I_{1}+I_{2})\cdot(R_{0}+r)$，

解得 $I_{1}=-\frac{R_{0}+r}{r_{1}+R_{1}+R_{0}+r}I_{2}+\frac{E}{r_{1}+R_{1}+R_{0}+r}$。

设图象的纵截距为 $b$、斜率为 $k$，由此式可知 $b=\frac{E}{r_{1}+R_{1}+R_{0}+r}$，$|k|=\frac{R_{0}+r}{r_{1}+R_{1}+R_{0}+r}$。延长图线与 $I_{1}$ 轴交点读出 $b$ 值，算出 $k$ 值，代入数据解得电源电动势 $E=1.4 \UV$，内阻 $r=0.5 \UO$。
}

\item
%% number: 10
%% paperName: 2017年天津理综第 10 题 16 分
%% typeId: 6
%% type: 计算题
%% chapter: 运动和力
%% point: 动量守恒定律
%% method: 无
%% score: 16
%% degree: 650
%% duplicateId: 0
%% body:
如图所示，物块 A 和 B 通过一根轻质不可伸长的细绳连接，跨放在质量不计的光滑定滑轮两侧，质量分别为 $m_{A}=2 \Ukg$、$m_{B}=1 \Ukg$。初始时 A 静止与水平地面上，B 悬于空中。先将 B 竖直向上再举高 $h=1.8 \Um$（未触及滑轮）然后由静止释放。一段时间后细绳绷直，A、B 以大小相等的速度一起运动，之后 B 恰好可以和地面接触。取 $g=10 \Umsq$。

\begin{enumerate}
	\item
	B 从释放到细绳绷直时的运动时间 $t$；
	\item
	A 的最大速度 $v$ 的大小；
	\item
	初始时 B 离地面的高度 $H$。
\end{enumerate}
\onepicture[h=4.4cm, align=right]{\includesvg{figs/10}}

%% answer: 
\jdanswer{
\begin{enumerate}
	\item
	$t=0.6 \Us$

	\item
	$v=2 \Ums$

	\item
	$H=0.6 \Um$

\end{enumerate}
}
%% memo:
\memoanswer{
（1）B 从释放到细绳刚绷直前做自由落体运动，有

$h=\frac{1}{2}gt^{2}$ ①，代入数据解得 $t=0.6 \Us$ ②。

（2）设细绳绷直前瞬间 B 速度大小为 $v_{B}$，有

$v_{B}=gt$ ③，

细绳绷直瞬间，细绳张力远大于 A、B 的重力，A、B 相互作用，由动量守恒得

$m_{B}v_{B}=(m_{A}+m_{B})v$ ④，

之后 A 做匀减速运动，所以细绳绷直后瞬间的速度 $v$ 即为最大速度，联立 ②③④ 式，代入数据解得

$v=2 \Ums$ ⑤。

（3）细绳绷直后，A、B 一起运动，B 恰好可以和地面接触，说明此时 A、B 的速度为零，这一过程中 A、B 组成的系统机械能守恒得

$\frac{1}{2}(m_{A}+m_{B})v^{2}+m_{B}gH=m_{A}gH$ ⑥，

代入数据解得 $H=0.6 \Um$ ⑦。
}

\item
%% number: 11
%% paperName: 2017年天津理综第 11 题 18 分
%% typeId: 6
%% type: 计算题
%% chapter: 磁场
%% point: 带电粒子在磁场中的运动
%% method: 无
%% score: 18
%% degree: 450
%% duplicateId: 0
%% body:
平面直角坐标系 $xOy$ 中，第Ⅰ象限存在垂直于平面向里的匀强磁场，第Ⅲ象限存在沿 $y$ 轴负方向的匀强电场，如图所示。一带负电的粒子从电场中的 Q 点以速度 $v_{0}$ 沿 $x$ 轴正方向开始运动，Q 点到 $y$ 轴的距离为到 $x$ 轴距离的 2 倍。粒子从坐标原点 $O$ 离开电场进入磁场，最终从 $x$ 轴上的 P 点射出磁场，P 点到 $y$ 轴距离与 Q 点到 $y$ 轴距离相等。不计粒子重力，问：

\begin{enumerate}
	\item
	粒子到达 $O$ 点时速度的大小和方向；
	\item
	电场强度和磁感应强度的大小之比。
\end{enumerate}
\onepicture[h=4.0cm, align=right]{figs/11.png}

%% answer: 
\jdanswer{
\begin{enumerate}
	\item
	$v=\sqrt{2}v_{0}$，方向与 $x$ 轴正方向的夹角为 $45^{\circ}$ 斜向上

	\item
	$\frac{E}{B}=\frac{v_{0}}{2}$

\end{enumerate}
}
%% memo:
\memoanswer{
（1）在电场中，粒子做类平抛运动，设 Q 点到 x 轴距离为 $L$，到 y 轴距离为 $2L$，粒子的加速度为 $a$，运动时间为 $t$，有

$2L=v_{0}t$ ①，

$L=\frac{1}{2}at^{2}$ ②，

设粒子到达 O 点时沿 y 轴方向的分速度为 $v_{y}$，

$v_{y}=at$ ③，

设粒子到达 O 点时速度方向与 x 轴正方向夹角为 $\alpha$，有

$\tan\alpha=\frac{v_{y}}{v_{0}}$ ④，

联立 ①②③④ 式得

$\alpha=45^{\circ}$ ⑤，

即粒子到达 O 点时速度方向与 x 轴正方向成 $45^{\circ}$ 角斜向上。设粒子到达 O 点时速度大小为 $v$，由运动的合成得

$v=\sqrt{v_{0}^{2}+v_{y}^{2}}$ ⑥，

联立 ①②③⑥ 式得

$v=\sqrt{2}v_{0}$ ⑦。

（2）设电场强度为 $E$，粒子电荷量为 $q$，质量为 $m$，粒子在电场中受到的电场力为 $F$，由牛顿第二定律得

$F=ma$ ⑧，

又 $F=qE$ ⑨，

设磁场的磁感应强度大小为 $B$，粒子在磁场中做匀速圆周运动的半径为 $R$，所受的洛伦兹力提供向心力，有

$qvB=m\frac{v^{2}}{R}$ ⑩，

由几何关系可知

$R=\sqrt{2}L$ \Circled{11}，

\onepicture[w=7cm]{figs/11_an.jpg}

联立 ①②⑦⑧⑨⑩\Circled{11} 式得

$\frac{E}{B}=\frac{v_{0}}{2}$ \Circled{12}。
}

\item
%% number: 12
%% paperName: 2017年天津理综第 12 题 20 分
%% typeId: 6
%% type: 计算题
%% chapter: 磁场
%% point: 安培力
%% method: 无
%% score: 20
%% degree: 450
%% duplicateId: 0
%% body:
电磁轨道炮利用电流和磁场的作用使炮弹获得超高速度，其原理可用来研制新武器和航天运载器。电磁轨道炮示意如图，图中直流电源电动势为 $E$，电容器的电容为 $C$。两根固定于水平面内的光滑平行金属导轨间距为 $l$，电阻不计。炮弹可视为一质量为 $m$、电阻为 $R$ 的金属棒 MN，垂直放在两导轨间处于静止状态，并与导轨良好接触。首先开关 S 接 1，使电容器完全充电。然后将 S 接至 2，导轨间存在垂直于导轨平面、磁感应强度大小为 $B$ 的匀强磁场（图中未画出），MN 开始向右加速运动。当 MN 上的感应电动势与电容器两极板间的电压相等时，回路中电流为零，MN 达到最大速度，之后离开导轨。问：

\begin{enumerate}
	\item
	磁场的方向；
	\item
	MN 刚开始运动时加速度 $a$ 的大小；
	\item
	MN 离开导轨后电容器上剩余的电荷量 $Q$ 是多少。
\end{enumerate}
\onepicture[w=11cm, align=right]{figs/12.png}

%% answer: 
\jdanswer{
\begin{enumerate}
	\item
	磁场的方向垂直于导轨平面向下

	\item
	$a=\frac{BlE}{mR}$

	\item
	$Q=\frac{B^{2}l^{2}C^{2}E}{m+B^{2}l^{2}C}$

\end{enumerate}
}
%% memo:
\memoanswer{
（1）垂直于导轨平面向下。

（2）电容器完全充电后，两极板间电压为 $E$，当开关 S 接 2 时，电容器放电，设刚放电时流经 MN 的电流为 $I$，有

$I=\frac{E}{R}$ ①，

设 MN 受到的安培力为 $F$，有

$F=BIl$ ②，

由牛顿第二定律得

$F=ma$ ③，

联立 ①②③ 式得

$a=\frac{BlE}{mR}$ ④。

（3）当电容器充电完毕时，设电容器上电量为 $Q_{0}$，有

$Q_{0}=CE$ ⑤，

开关 S 接 2 后，MN 开始向右加速运动，速度达到最大值 $v_{\max}$ 时，设 MN 上的感应电动势为 $E'$，有

$E'=Blv_{\max}$ ⑥，

依题意有

$E'=\frac{Q}{C}$ ⑦，

设在此过程中 MN 的平均电流为 $\bar{I}$，MN 上受到的平均安培力为 $\bar{F}$，有

$\bar{F}=B\bar{I}l$ ⑧，

由动量定理，有

$\bar{F}\Delta t=mv_{\max}-0$ ⑨，

又 $\bar{I}\Delta t=Q_{0}-Q$ ⑩，

联立 ⑤⑥⑦⑧⑨⑩ 式得

$Q=\frac{B^{2}l^{2}C^{2}E}{m+B^{2}l^{2}C}$ \Circled{11}。
}

\end{enumerate}

\end{document}
